%%%%%
%
\section{Vent apparent: vecteur unique fait � partir de deux vecteurs}

\itemstyle{Op\'erateur~:} \operfmt{choix\_vent\_app}
\index{choix\_vent\_app}

%%%%%OBJET%%%%
\itemstyle{Objet~:}
Choix d'une mesure de vent apparent � partir de 2 capteurs
de m�me mod�le situ�s de 2 cot� d'un m�t de mesure.\newline
  Les crit�res de choix sont relativement empiriques.


%%%%%ARGUMENTS%%%%
\bigskip
%%%%%ENTREES%%%%
\begin{tabarg}
\multicolumn{3}{l}{\titretab {Entr\'ees}}\\
\hline
$dd1$ & Vecteur, {\tt float} & Direction vent apparent babord ($^\circ$)\\
$ff1$ & Vecteur, {\tt float} & Force vent apparent babord ($kt$, $0$ 
correspondant � un vent apparent venant de l'avant du bateau)\\
$dd2$ & Vecteur, {\tt float} & Direction vent apparent tribord ($^\circ$)\\
$ff2$ & Vecteur, {\tt float} & Force vent apparent tribord ($kt$, $0$
correspondant � un vent apparent venant de l'avant du bateau)\\
\end{tabarg}

\medskip
%%%%%SORTIES%%%%
\begin{tabarg}
\multicolumn{3}{l}{\titretab {Sorties}}\\
\hline
$dd\_ap$ & Vecteur, {\tt float} & Direction vent apparent ($^\circ$)\\
$ff\_ap$ & Vecteur, {\tt float} & Force vent apparent ($kt$)\\
\end{tabarg}


%%%%%FORMULE%%%%
\itemstyle{Formule}

%%choix_vent_app
   
   - on choisit l'ensemble an�mo-girouette qui est au vent, c'est � dire:\newline
     dd1-ff1 �tant � babord et dd2-ff2 � tribord,\newline
     si dd1 $>$ 0 on choisit la girouette 2\newline
     si dd1 $<$ 0 on choisit la girouette 1\newline
   - par vent de face, on choisit l'ensemble an�mo-girouette qui mesure la force 
   du vent la plus �lev�e.


%%%% COMMENTAIRE %%%%
\itemstyle{Commentaire~:}
   
Lorsque les girouettes sont mont�es avec le $0$ pointant vers
l'arri�re du bateau (cas d'EQUALANT, POMME), il faut que 
la rotation corrective soit faite \textbf{avant} d'appliquer
cet op�rateur.

\textbf{ATTENTION} : cet op�rateur contient un test sp�cifique aux
an�mom�tres YOUNG utilis�s pendant EQUALANT-99 et POMME : On �limine la
girouette 2 si la direction est sup�rieure � $171^\circ$ 
ou inf�rieure � $178^\circ$, (apr�s �tude des 
histogrammes de zones mortes des girouettes).
   
%%%%%AUTEUR%%%%
\itemstyle{Auteur~:}
CNRM/GMEI/TRAMM

%%%%%REFERENCE%%%%
\itemstyle{R\'ef\'erences~:}
Campagne pr�c�dante FETCH
